\chapter{Conclusiones Finales}
\label{ch:conclusiones}

El análisis sobre los determinantes de la solvencia intertemporal de la deuda consolidada del SPNF argentino (2004--2025) permite arribar a síntesis conceptuales que trascienden la exégesis econométrica y abordan la economía política del endeudamiento periférico, manteniendo el rigor estadístico para distinguir evidencia de conjeturas.

\section{Síntesis de Hallazgos Empíricos}

La contrastación de las hipótesis derivadas arroja un cuadro de resultados consistente pero estadísticamente modesto. La Tabla \ref{tab:mapa_consistencia} resume la correspondencia entre cada hipótesis, las variables y el método empleados para contrastarla, el resultado obtenido y el veredicto que de él se sigue; los párrafos siguientes desarrollan cada caso.

\begin{table}[H]
\centering
\caption{\label{tab:mapa_consistencia} Mapa de Consistencia: Hipótesis, Método, Resultado y Veredicto.}
\footnotesize
\begin{tabular}{p{2.6cm}p{2.9cm}p{2.7cm}p{3.3cm}p{3.5cm}}
\toprule
\textbf{Hipótesis} & \textbf{Variables} & \textbf{Método} & \textbf{Resultado} & \textbf{Veredicto} \\
\midrule
$H_{1}$ --- Reacción fiscal débil & $pb_t$, $d_{t-1}$, $\tilde{y}_t$, EMBI+, TCRM & DOLS ($m=4$, HAC, ventana original) y VECM (ventana ampliada y original) & DOLS: $\rho=-0.0071$ ($p=0.564$). VECM ampliada: $\alpha_{pb}=0.0044$ ($p=0.204$). VECM original: $\alpha_{pb}=0.0014$ ($p=0.559$) & \textbf{No rechazada} en las tres estimaciones: sin reacción sistemática \\
\midrule
$H_{1a}$ --- Fatiga fiscal (umbral no lineal) & $pb_t$, $d_{t-1}$, EMBI+ como variable de umbral & Umbral de Hansen, Sup-LM por remuestreo estocástico (1.000 rép.) & Niveles: $\tau^*=449$ pb, Sup-LM $p=0.076$ (partición 14 vs. 74 obs.). $\Delta pb_t$ (I(0), Sec. \ref{sec:robustez_hansen_dpb}): $\tau^*=2083$ pb, Sup-LM $p<0.001$ & \textbf{Evidencia mixta pero ya no unilateral}: marginalmente significativa en niveles (partición extrema, coeficientes individuales no significativos); sólidamente significativa en $\Delta pb_t$ \\
\midrule
$H_{1b}$ --- Endogeneidad del riesgo soberano & $pb_t$, EMBI+ instrumentado con VIX y diferencial regional & IV-2SLS; Wu-Hausman, Sargan, F de 1.\textsuperscript{a} etapa & Endogeneidad confirmada ($p=0.009$); instrumentos fuertes ($F=13.00$) y válidos (Sargan $p=0.436$); coeficiente $=0.0017$ ($p=0.076$) & \textbf{Parcial, con señal marginal}: endogeneidad sí, canal causal débil pero ya no descartable \\
\midrule
$H_{1c}$ --- Quiebres estructurales & $d_t$, $d_t$ consolidada & Zivot-Andrews (quiebre único); Bai-Perron (múltiples, BIC) & Z-A no significativo ($t=-2.83$); Bai-Perron: 3 quiebres (SPNF) y 2 (consolidada) & \textbf{Verificada} en su versión de quiebres múltiples \\
\midrule
Extensión --- Solvencia prospectiva & $d_t$, $pb_t$, $g_t$, $r_t$, $\Delta e_t$ & DSA Monte Carlo, $t$ de Student ($\nu\approx 4.8$, $\lambda\approx 0.076$), GARCH/DCC & $P(d_{2035}>100\%)=31.2\%$ (32.5\% con DCC-GARCH) & \textbf{Riesgo de cola relevante}, no desenlace modal \\
\bottomrule
\end{tabular}
\begin{flushleft}
\footnotesize \textit{Nota.} Elaboración propia. ``No rechazada'' se refiere a la hipótesis nula econométrica de ausencia de reacción, que es la hipótesis sustantiva $H_1$ de esta tesis. VECM ampliada: ventana 1999--2025 ($n=108$), rango de cointegración impuesto por teoría (Sección \ref{sec:resultados_vecm}). VECM original: ventana 2004--2025 ($n=87$), rango no impuesto, emerge del contraste de Johansen. Las secciones de origen de cada resultado se detallan en los párrafos siguientes y en el Capítulo \ref{ch:resultados}.
\end{flushleft}
\end{table}

Respecto de $H_1$ (reacción fiscal), tres estimaciones independientes convergen en el mismo veredicto. La DOLS de la Función de Reacción Fiscal no permite rechazar la hipótesis nula de ausencia de reacción sistemática del superávit primario frente al endeudamiento heredado ($\rho=-0.0071$, $p=0.564$). El VECM que la complementa como técnica de referencia sobre la ventana ampliada, a pedido del director de tesis (Sección \ref{sec:resultados_vecm} del Capítulo \ref{ch:resultados}), trata las cuatro variables del sistema como conjuntamente endógenas en lugar de purgar la endogeneidad con una construcción de ecuación única, y llega al mismo resultado nulo ($\alpha_{pb}=0.0044$, $p=0.204$); lo mismo ocurre en el VECM estimado sobre la ventana original con series reales, donde el rango de cointegración ya no necesita imponerse por teoría sino que emerge limpiamente del propio contraste de Johansen ($\alpha_{pb}=0.0014$, $p=0.559$). Que tres estrategias de estimación con supuestos de identificación tan distintos coincidan en el mismo diagnóstico es, para esta tesis, la evidencia más robusta a favor de $H_1$: no es un artefacto de una técnica en particular. La ampliación de la ventana muestral, lejos de revertir el resultado nulo, que era una posibilidad real dado el pedido explícito del director de extender la muestra ``aun con cambios estructurales extremos'', lo confirma sobre datos adicionales que incluyen el colapso de la convertibilidad y el default de 2001--2002.

Respecto de $H_{1a}$ (fatiga fiscal y umbrales no lineales), el modelo de umbral de Hansen, en su especificación en niveles, la que replica directamente la forma canónica de \citet{ghosh2013}, localiza sobre la serie real de EMBI+ un quiebre candidato en 449 puntos básicos y un ordenamiento de coeficientes coincidente con el que predeciría ese modelo (0.0305 en régimen ``normal'' frente a 0.0126 en régimen de estrés); a diferencia del diagnóstico obtenido con la serie de contingencia de EMBI+ empleada hasta la Sección \ref{sec:auditoria_fallback} del Capítulo \ref{ch:metodologia}, el contraste por remuestreo estocástico Sup-LM del propio umbral ahora sí rechaza la linealidad, aunque solo al 10\% ($p=0.076$), y sobre una partición muestral extremadamente desigual (14 contra 74 observaciones) que explica por qué ninguno de los dos coeficientes de régimen es individualmente significativo pese a ese rechazo global.

Esta especificación en niveles, además, viola el supuesto de estacionariedad estricta que el propio test de Hansen exige de su variable dependiente, dado que $pb_t$ es I(1) (Sección \ref{sec:dfgls}). La Sección \ref{sec:robustez_hansen_dpb} reestima el contraste con $\Delta pb_t$, estacionaria y por tanto la única de las dos especificaciones que satisface formalmente ese supuesto, y el resultado es contundente: el contraste por remuestreo estocástico Sup-LM rechaza la linealidad con holgura ($p<0.001$) y ambos coeficientes de régimen son significativos, con el patrón exacto de atenuación bajo estrés que predice la fatiga fiscal, sobre una partición muestral (74 contra 14, en sentido inverso a la de niveles) que no concentra el resultado en un puñado de observaciones. El veredicto de esta tesis para $H_{1a}$ es, en consecuencia, de \textbf{evidencia mixta pero ya no unilateralmente débil}: la especificación formalmente inválida (niveles) ofrece una señal frágil pero presente; la formalmente válida ($\Delta pb_t$) ofrece la evidencia más sólida de umbral de toda la tesis, aunque para un objeto, la reacción del \textit{cambio} en el resultado primario, distinto de su \textit{nivel}, no idéntico al que \citet{ghosh2013} modela. Esta tesis no interpreta ninguno de los dos resultados de manera aislada como confirmación o refutación definitiva; el Capítulo \ref{ch:discusion} desarrolla la tensión entre ambos, y explicita que la lectura de economía política que se les asocia es una interpretación compatible con la evidencia, no una inferencia derivada unívocamente de ella.

Respecto de $H_{1b}$ (endogeneidad del EMBI+ sobre la reacción fiscal), la estimación IV-2SLS, instrumentando el EMBI+ con el VIX y un diferencial de riesgo soberano regional (diferencial de mercado sobre el ETF EMB, en reemplazo del TCRM rezagado de una especificación preliminar), produce instrumentos fuertes en primera etapa (F$=13.00$, por encima del umbral de 10 de \textcite{staiger1997}, aunque con menor margen que con la serie de contingencia de EMBI+) y válidos en conjunto (Sargan $=0.608$, $p=0.436$); el test de Wu-Hausman confirma, además, la endogeneidad del riesgo soberano ($7.32$, $p=0.009$). Con este diseño instrumental válido, el coeficiente del EMBI+ instrumentado ($0.0017$) alcanza significatividad marginal al 10\% ($p=0.076$): una señal débil, sobre una muestra reducida a $n=73$ por la cobertura del instrumento, pero que ya no permite sostener sin matices la ausencia de canal que sugería la serie de contingencia ($p=0.297$). Esta tesis no fuerza esta señal marginal hacia una conclusión firme; la reporta con la misma cautela con que reporta el resto de la evidencia mixta del capítulo.

Respecto de $H_{1c}$ (quiebres estructurales históricos), el contraste de Zivot-Andrews con quiebre endógeno localiza el punto de ruptura de la ratio Deuda/PIB en el segundo trimestre de 2018, coincidente con el cierre del mercado voluntario de crédito y la aprobación del acuerdo \textit{stand-by} con el FMI, pero el estadístico de prueba ($t=-2.83$) no alcanza significatividad ni siquiera al 10\% frente a los valores críticos de \textcite{zivot1992}. La ratio de deuda argentina se comporta, en términos estrictamente estadísticos, como una serie genuinamente integrada de orden uno, sin una tendencia determinística estable ni siquiera reconociendo el quiebre más evidente de la serie. La restricción de Zivot-Andrews a un único quiebre se supera mediante el procedimiento formal de \textcite{bai2003} (programación dinámica exacta, selección del número de quiebres por BIC), que identifica $H_{1c}$ con mayor riqueza: tres quiebres en la deuda SPNF (2007T2, 2014T3 y 2018T1, este último muy próximo al candidato de Zivot-Andrews, un trimestre antes) y dos en su versión consolidada ampliada con los pasivos remunerados del BCRA (2007T2 y 2016T4). A diferencia del test de quiebre único, $H_{1c}$ encuentra aquí soporte estadístico afirmativo: la serie sí exhibe quiebres estructurales identificables, aunque múltiples y no reducibles a uno solo.

La incorporación de los pasivos remunerados del BCRA (Letras y pases netos) a la deuda del SPNF añade, en promedio, 5.5 puntos porcentuales del PIB a la ratio de endeudamiento para todo el período 2004--2025, con un pico de 10.6\% en el primer trimestre de 2018 y una convergencia a valores prácticamente nulos desde mediados de 2024, cuando el Tesoro asumió el rol de esterilizador monetario mediante Letras Fiscales de Liquidez (LEFI). Esta capa cuasi-fiscal, ausente de la definición central de $d_t$ por decisión metodológica (Capítulo \ref{ch:metodologia}), modifica la cronología de quiebres estructurales sin alterar el diagnóstico central de $H_1$: la deuda consolidada ampliada (SPNF + BCRA) conserva el mismo orden de integración I(1) que la deuda consolidada del SPNF.

Finalmente, la proyección estocástica de sostenibilidad (DSA, Monte Carlo con perturbaciones $t$ de Student, $\nu\approx 4.8$, $\lambda\approx 0.076$, 1.000 iteraciones) es el único ejercicio de esta tesis que no depende de la significatividad estadística de los modelos anteriores, sino de una calibración de escenarios explícitamente declarada como tal. Arroja una probabilidad del 31.2\% de que la ratio Deuda/PIB supere la frontera del 100\% del producto hacia 2035, con cambios moderados (32.5\% con DCC-GARCH) al sustituir la calibración original por desvíos GARCH(1,1) y la correlación condicional dinámica de un modelo DCC-GARCH (Capítulo \ref{ch:discusion}). Se trata de una probabilidad de cola relevante para el análisis macroprudencial, aunque de ningún modo un desenlace mayoritario, y robusta a la sustitución de la calibración por estimación directa.

\section{Aporte Metodológico y Límites del Estudio}
\label{sec:aporte_metodologico}

El aporte de esta tesis reside en someter la hipótesis de fatiga fiscal a un protocolo econométrico secuencial completo y en reportar el resultado que ese protocolo arroja, predominantemente negativo o mixto según el objeto exacto contrastado. La contribución no es la confirmación de una tesis previa sino la delimitación precisa de lo que la evidencia argentina 2004--2025 permite y no permite sostener: no permite sostener una regla de reacción fiscal activa ni un canal causal significativo entre riesgo soberano y esfuerzo primario; permite sostener un umbral de fatiga fiscal solo de manera parcial y sensible a la especificación ---no significativo en niveles, significativo en primera diferencia (Sección \ref{sec:robustez_hansen_dpb})---; y sí permite sostener que la ratio de deuda exhibe quiebres estructurales múltiples identificables, que el parámetro de reacción es inestable entre subperíodos, y que la probabilidad de superar el 100\% del PIB hacia 2035 constituye un riesgo de cola no despreciable bajo supuestos explícitos.

De ello se sigue el desplazamiento del eje explicativo que articula esta tesis: si el mecanismo de ajuste no transita por la función de reacción fiscal del Tesoro (regla de Bohn, Nivel III del marco teórico), la explicación debe buscarse en el canal que el Paradigma de la Restricción Externa identifica (Nivel I): tipo de cambio real, efecto de hoja de balance y Ajuste Stock-Flujo, tal como desarrolla el Capítulo \ref{ch:discusion}. Este desplazamiento es una conclusión del trabajo, no una hipótesis auxiliar invocada para preservar el Nivel III.

Entre las limitaciones explícitamente reconocidas se cuentan: (i) el tamaño muestral moderado (88 observaciones trimestrales, reducido a 73 en la especificación IV-2SLS por la cobertura del diferencial regional de riesgo desde 2007), que reduce el poder estadístico de las técnicas de umbral y de instrumentación; (ii) la autocorrelación residual en la estimación DOLS (Durbin-Watson $=0.321$), reflejo de la inercia estructural del proceso presupuestario argentino más que de una falla de ajuste, aunque corregida en la inferencia mediante errores HAC; (iii) la persistencia de la calibración, en lugar de la estimación directa, para dos de los cinco impulsos o perturbaciones de la matriz de covarianzas del DSA (las tasas de interés doméstica e internacional, sin serie histórica propia en la base de datos construida), pese a que los tres impulsos restantes ya se reestiman vía GARCH(1,1) y DCC-GARCH (Capítulo \ref{ch:discusion}); (iv) la heterocedasticidad condicional confirmada mediante el test ARCH-LM sobre los residuos del DOLS ($LM=18.74$, $p=0.0009$, Sección \ref{sec:dfgls}); corregida en la inferencia mediante errores HAC, pero indica que la varianza de las perturbaciones sobre el resultado primario no es constante en el tiempo, un patrón que una especificación lineal estática no modela explícitamente; (v) la inestabilidad paramétrica de la Función de Reacción Fiscal a lo largo de la muestra (Sección \ref{sec:sensibilidad_temporal}): el coeficiente $\rho$ es positivo y significativo al 5\% o al 10\% en tres de las cuatro especificaciones DOLS por subperíodo, pero nulo en la estimación de muestra completa, lo que impide interpretar el coeficiente agregado como el parámetro de una regla de comportamiento estable y advierte sobre la sensibilidad de cualquier estimación de la regla de Bohn a la ventana temporal y a la especificación dinámica elegidas; y (vi) la ausencia de evidencia de cointegración bivariada mediante el test residual de Engle-Granger entre deuda y resultado primario (Apéndice \ref{sec:engle_granger}, $p=0.3012$), consistente con el diagnóstico divergente de Johansen y que refuerza la cautela con que deben interpretarse las estimaciones de largo plazo de este capítulo.

Su jerarquización orienta la agenda futura. La más \textbf{estructural} es la (i): el tamaño muestral de un único país con 88 observaciones trimestrales es el límite fundamental que reduce la potencia de todas las técnicas de umbral, cointegración e instrumentación aplicadas. No es superable dentro del marco de serie de tiempo de un único país salvo esperando que transcurra el tiempo, pero sí lo es mediante la extensión a un panel de países comparables (Sección \ref{sec:agenda_futura}). La más \textbf{superable} a corto plazo es la (iii): la cobertura del diferencial de riesgo soberano regional (ETF EMB) desde 2007 es una limitación de fuente, no de diseño, que podría resolverse con acceso a bases de datos de pago (JP Morgan EMBI Historical, Bloomberg) o con la construcción de una variable de contraste regional alternativa sobre el período 2004--2007, recuperando las 14 observaciones excluidas de la estimación IV-2SLS y aumentando su potencia estadísticamente relevante. Las restantes limitaciones (ii), (iv), (v) y (vi) son en buena medida inherentes a la naturaleza de la economía argentina (su volatilidad estructural, su inercia fiscal y sus quiebres de régimen frecuentes) y no constituyen fallas de diseño superables con un rediseño alternativo dentro de los mismos datos.

\section{Agenda de Investigación Futura e Implicancias Analíticas}
\label{sec:agenda_futura}

El horizonte posterior a 2025 demanda combinar la prudencia estadística de esta tesis con el análisis riguroso de las implicancias macroeconómicas que sus hallazgos, aun modestos, permiten sostener. Las proyecciones de organismos multilaterales que asumen superávits primarios permanentes de 3\% o 4\% del PIB sin contemplar la ausencia estadísticamente verificada de una reacción fiscal sistemática de esa magnitud en el historial argentino reciente deberían, como mínimo, ponderar el escenario de Estrés y la probabilidad de cola del 31.2\% documentados en el Capítulo \ref{ch:discusion}.

En términos analíticos, el hallazgo empírico de una correlación positiva, aunque débil, entre las variaciones del tipo de cambio real y la dinámica de la ratio de deuda ($r=0.14$ sobre la serie real de TCRM, Sección \ref{sec:hoja_balance} del Capítulo \ref{ch:descriptivo}) señala que la exposición al riesgo cambiario y la estructura de plazos de los pasivos constituyen determinantes moderados, no dominantes en la correlación bivariada simple, del Ajuste Stock-Flujo. Ello implica que la modelización econométrica de la solvencia en economías emergentes debe incorporar de forma explícita las perturbaciones del tipo de cambio real y la composición por monedas de los pasivos, en lugar de confiar exclusivamente en metas de resultado primario cuyo carácter estabilizador no encuentra sustento estadístico en la evidencia reciente.

Esta tesis incorporó cinco extensiones metodológicas respecto de su diseño original: la estimación VECM sobre una ventana muestral ampliada mediante empalme histórico (1999--2025, $n=108$), que reemplaza a DOLS como técnica de referencia a pedido del director de tesis (Capítulos \ref{ch:metodologia}--\ref{ch:resultados}); el procedimiento formal de quiebres múltiples de \textcite{bai2003} (Capítulo \ref{ch:descriptivo}); la sustitución del TCRM rezagado por un diferencial de riesgo soberano regional como instrumento del EMBI+ (Capítulo \ref{ch:resultados}); la reestimación GARCH(1,1)/DCC-GARCH de la matriz de covarianzas del DSA (Capítulo \ref{ch:discusion}); y la consolidación de la deuda SPNF con los pasivos remunerados del BCRA (Capítulos \ref{ch:metodologia}--\ref{ch:descriptivo}). En términos de agenda de investigación futura, esta tesis identifica cinco líneas de continuidad directa.

\textbf{Primera}, y con prioridad sobre las restantes por seguirse directamente del resultado mixto y sensible a la especificación del contraste de umbral (Sección \ref{sec:robustez_hansen_dpb}), contrastar la no linealidad de la Función de Reacción Fiscal mediante especificaciones de \textbf{transición suave} (STAR/LSTR) y de \textbf{umbrales múltiples}. La formulación de \citet{hansen1999} supone una transición discreta entre dos regímenes; si el ajuste de la conducta fiscal frente a la presión financiera es gradual y no abrupto, como sugiere la discusión del Capítulo \ref{ch:discusion}, esa clase de modelo es la especificación adecuada para detectarlo, y su rechazo o no rechazo constituiría un contraste genuino de la hipótesis de fatiga fiscal, no una reinterpretación del resultado nulo aquí obtenido. Convendría además evaluar variables de umbral alternativas al EMBI+ (ratio de deuda en moneda extranjera, cobertura de reservas sobre vencimientos anuales), cuya menor volatilidad relativa mejoraría la relación señal-ruido del contraste.

\textbf{Segunda}, implementar el contraste secuencial \textit{sup} F($l$+1$\vert$$l$) de Bai y Perron con sus tablas de valores críticos no estándar, que permitiría construir intervalos de confianza para la ubicación de cada quiebre, allí donde el criterio BIC aquí empleado solo informa el número óptimo de quiebres.

\textbf{Tercera}, extender la cobertura del diferencial soberano regional a 2004--2007 (sin fuente gratuita disponible en este trabajo) para recuperar las 14 observaciones excluidas de la estimación IV-2SLS, y explorar instrumentos adicionales (medidas de riesgo político-electoral) que permitan contrastar la robustez del hallazgo de no significatividad del canal EMBI+.

\textbf{Cuarta}, retomar el análisis de dominancia fiscal de \textcite{sargent1981} aprovechando la serie de pasivos remunerados del BCRA ya construida (Sección \ref{sec:deuda_consolidada}), que esta tesis empleó para consolidar el stock de deuda pero no para estimar formalmente un régimen de dominancia monetaria-fiscal.

\textbf{Quinta}, y con mayor potencial transformador para la agenda, extender el análisis a un panel de economías emergentes latinoamericanas comparables (Brasil, Uruguay, Colombia, Chile, Perú) utilizando datos del Fondo Monetario Internacional y la CEPAL para el período 1990--2025. Un panel de 6 a 8 países con series de 35 años proveería una muestra de entre 700 y 1.000 observaciones anuales, aumentando la potencia estadística de las pruebas de umbral y de la Función de Reacción Fiscal en varios órdenes de magnitud respecto del caso de un único país. Este enfoque de panel permitiría, además, identificar si la debilidad estadística de la reacción fiscal que esta tesis documenta para Argentina es un fenómeno idiosincrático de su historia de cesaciones de pagos y regresiones políticas, o si refleja una condición más amplia de la economía política fiscal periférica en América Latina.

La ausencia de una regla de reacción fiscal verificable, documentada en esta tesis, es en sí misma un dato sobre la naturaleza de la solvencia soberana argentina contemporánea: una solvencia que, en el período estudiado, no se sostuvo sobre el ajuste primario sino sobre la licuación y la reestructuración de pasivos.

\section{Límites Cognitivos del Estudio y Consideraciones Prospectivas}

El objetivo general de esta tesis es una meta de conocimiento puro, no una propuesta de intervención de política económica. El descubrimiento de la inoperancia de la regla de Bohn y la evidencia mixta ---sensible a la especificación--- sobre el umbral de Hansen representan hallazgos cognitivos sobre las determinaciones materiales del vector fiscal argentino. Por ende, los hallazgos analíticos sobre la vulnerabilidad cambiaria y el perfil de vencimientos se presentan solo a título prospectivo como apertura heurística hacia futuras líneas de investigación. La transformación de un problema de conocimiento en una acción pragmática excede el alcance empírico de este trabajo.
